\section{Parameters and Primitives}

We define SHRINCS-SHA2-128s as a SHRINCS instance preferable for Bitcoin. It operates with the 
following set of parameters:

\begin{center}
\begin{tabular}{ c c c }
 Parameter & Value & Description \\ 
 \hline
 \texttt{H} & \texttt{SHA-256} & The underlying hash function \\ 
 \texttt{n} & $128$ & Security parameter (bits) \\  
 \texttt{w} & $256$ & Winternitz parameter (8 bits per chain) \\
 \texttt{l}  & $16$ & Message digest length, the number of WOTS+C chains \\
 \texttt{Swn} & $2040$ & Target sum for WOTS+C \\
 \texttt{hsf} & $\color{red}10\color{black}$ & Maximum stateful tree height \\
 \texttt{hsl} & $\color{red}63\color{black}$ & Maximum stateless hypertree height \\
 \texttt{d} & $\color{red}7\color{black}$ & Stateless hypertree layers \\
 \texttt{a} & $\color{red}12\color{black}$ & FORS/PORS tree height \\
 \texttt{k} & $\color{red}14\color{black}$ & Number of FORS/PORS trees \\
 \texttt{c} & $4$ & Counter size for WOTS+C grinding (bytes) \\ 
\end{tabular}
\end{center}

The foundation of SHRINCS is the hash function SHA-256, operated in specific modes to construct 
WOTS+C (Compressed Winternitz One-Time Signatures) and Unbalanced Merkle Trees. At the same time
SHRINCS utilizes the SLH-DSA for a stateless signature instance. 

\subsection{Hash Function Instantiation}

SHRINCS uses SHA-256 truncated to \texttt{n} bits for all internal hash operations. 
\begin{center}
\texttt{H(m) = SHA-256(m)[0:n]}
\end{center}

\subsection{Tweakable Hash Function}

To prevent multi-target attacks and ensure domain separation between the various components 
(WOTS+, trees, stateless fallback, etc), we employ the Tweakable Hash Function construction defined in 
FIPS 205:

\begin{center}
\texttt{Tw(P, ADRS, m) = H(P || ADRS || m)}
\end{center}
where: 
\begin{itemize}
    \item[] \texttt{P} is a public parameters
    \item[] \texttt{ADRS} is a tweak  (context address)
    \item[] \texttt{m} is a message
\end{itemize}

This usage of tweaked hashing is critical. It ensures that a hash computed at one position in the 
Merkle tree cannot be substituted for a hash at another position, effectively mitigating generic 
tree-grafting attacks.

\subsubsection{Tweak Structure}

Tweaks are structured as follows to ensure uniqueness across all scheme components:

\begin{center}
\texttt{ADRS = layer || tree\_address || type || words}
\end{center}
where: 
\begin{itemize}
    \item[] \texttt{layer}: (word 0)layer in hypertree (0 = bottom)
    \item[] \texttt{tree\_address}: (words 1-3) tree address within layer
    \item[] \texttt{type}: (word 4) address type (Appendix~A)
    \item[] \texttt{position}: (words 5-7) type-dependent 12 bytes (Appendix~B)
\end{itemize}

\subsection{Pseudorandom Function}
We use PRF for deterministic key derivation

\texttt{PRF(ADRS, SK.seed) = H(0x06 || SK.seed)[0:n]}

\subsection{Message Hash Function}

\texttt{H\_msg(ADRS, R, PK.seed, PK.root, m) = H(0x07 || R || PK.seed || PK.root || m)}
